\section{Discussion and Conclusions}

In this study, we have shown that a simple three‐population firing‐rate network model, endowed with a non‐selective urgency ramp and transient sensory inputs, can quantitatively reproduce the key behavioral signatures of mice performing the 3CDR task. As illustrated in Figure~\ref{fig:sim-trial-types}, the model captures the decline in choice accuracy from Easy through Hard trials, and—when applied to the Variable–Stimulus and Variable–Delay protocols—it likewise predicts performance drops for both very brief and very late cues (Figures~\ref{fig:sim-stim-dur} and \ref{fig:sim-delay-dur}). These findings suggest that errors arise primarily from the temporal misalignment between evidence availability and the rising urgency signal, rather than from an intrinsic decay of memory during the delay.

By performing a normal‐form reduction up to cubic order, we have further shown that the network dynamics can be visualized as a particle diffusing on an “energy landscape” that initially features a central well (the no‐decision state) and, as urgency increases, evolves into a landscape with three peripheral wells corresponding to each choice. This geometric picture (Figure~\ref{fig:potential-o3}) intuitively explains how early stimulus offset can leave the system trapped in the central basin, while late stimulus onset can cause commitment before evidence arrives (Figure~\ref{fig:timing_illustration}).

Our model generalizes urgency‐signal based frameworks developed for two‐choice decision-making \cite{inagaki2019} to the case of multiple alternatives in a principled way. Unlike classical DDMs that assume infinite memory, our formulation shows that even short delays can produce near‐chance performance if they misalign with an urgency signal, shifting the interpretation of “memory limitations” in delayed‐response tasks. Nonetheless, we have made simplifying assumptions: urgency is implemented as a fixed linear ramp and we have not yet fitted or modeled individual mice’s idiosyncratic biases or their trial‐by‐trial adaptation.

Future work will focus on three main extensions. First, we will endow each excitatory population with a synaptic plasticity rule for its self‐excitation coefficient, allowing the model to learn and generate biases based on trial history. Second, we plan to fit the reduced normal‐form model directly to individual mice’s psychometric and reaction‐time data, thereby calibrating urgency and stimulus slopes, amplitudes, onsets, offsets and general noise levels against empirical variability. Third, by exploiting the energy‐landscape framework, we will apply escape‐rate theory to compute the probability of transitions between attractor wells, yielding quantitative predictions for choice reversals as a function of stimulus timing and urgency dynamics.

In conclusion, by unifying an urgency signal framework, transient sensory evidence and attractor‐based competition within an explicit energy‐landscape formalism, our model provides a powerful quantitative tool for interpreting both behavioral and neural data in multi‐alternative decision‐making tasks, and generates clear, testable predictions for future perturbation studies.


This work was presented as a poster at the annual Barcelona Computational, Cognitive and Systems Neuroscience (BARCCSYN) meeting 2025 and was awarded “Best Poster” in Session 1.